\section{Detailed intervention results and numerical tables}
\label{app:detailed-results}

This appendix retains the full numerical comparisons behind the three main-text findings. It reports the region-specific behavior and hidden-position grids, the same-object sparse effects and descriptive summaries, native-channel path tests, and the shared-subspace and complete-state comparisons. The main text states the scientific result of each contrast; the tables here preserve all means, intervals, controls, and selection boundaries.

\subsection{Localized visual evidence has region-specific answer effects in all three models}
\label{sec:shared-dependence}

The first question is whether removing the annotated answer region lowers target-answer scores more than removing a comparable non-evidence region in all three models. A second requirement connects that regional loss to internal computation: restoring the aligned clean hidden state must recover more score under answer-region corruption than under matched random corruption. The first comparison establishes regional specificity, the second connects the regional loss to an internal state, and a frozen control set tests whether the recovered content, direction, and location are specific.

\subsubsection{The annotated region matters beyond generic occlusion}

Figure~\ref{fig:behavior-region-specificity} compares answer-region corruption with same-area random corruption on the same image--question examples. Positive values mean that the annotated region causes greater damage to the target rank or complete-answer score. In the expanded 61-example evaluation, the rank intervals lie above zero for all three models, and the complete-answer intervals lie above zero for Qwen and LLaVA; the remaining Gemma interval includes zero. Thus, five of the six model-by-metric estimates lie strictly above zero. In the independently selected 52-example evaluation, both outcomes have intervals above zero in all three models. Removing the annotated answer region therefore causes more damage than removing an equally sized arbitrary region.

\subsubsection{Hidden restoration connects the regional loss to internal computation}

Output damage alone does not show whether the target-score loss can be recovered from an internal model state. We therefore restore clean residual states during the answer-region-corrupted and same-area-random-corrupted runs. For each predefined position group, we compare the restoration benefit under answer-region corruption with that under matched same-area random corruption. The main comparison reports answer-adjacent positions and all visual positions pooled together, so neither position group is selected by improvement in the target answer.

Table~\ref{tab:hidden-position-grid} reports the 61-example comparison. Gemma and Qwen are clearest when all visual positions are restored together, whereas LLaVA is clearest at the predefined answer-adjacent positions. Under those model-specific position groups, both the target-rank and target-logit intervals lie above zero. For LLaVA, the all-visual intervals instead cross zero. Thus, the same region-specific answer-score pattern need not be recovered most strongly at the same type of position in every model.

\begin{center}
\begin{minipage}{\linewidth}
\captionsetup{type=table,hypcap=false}
\centering
\caption{\textbf{Hidden restoration recovers target-answer scores at different position groups.} Entries report the restoration benefit under answer-region corruption minus that under matched same-area random corruption on 61 examples. $\Delta s$ is the mean score change, $\pi$ the positive-sample fraction, and brackets the 95\% sample-bootstrap interval. AV denotes all visual positions pooled together and AA denotes answer-adjacent positions. Rank changes are interpreted within model because vocabulary sizes differ.}
\label{tab:hidden-position-grid}
\normalsize
\begin{tabular*}{\linewidth}{@{\extracolsep{\fill}}lcccc@{}}
\toprule
Model--position & $\Delta s_{\mathrm{rank}}\uparrow$ [95\% CI] & $\pi_{\mathrm{rank}}$ & $\Delta s_{\mathrm{logit}}\uparrow$ [95\% CI] & $\pi_{\mathrm{logit}}$ \\
\midrule
Gemma (AV) & $14{,}444.7\ [7{,}248.8,23{,}404.2]$ & .738 & $4.818\ [2.421,7.249]$ & .656 \\
Qwen (AV) & $2{,}711.5\ [549.4,6{,}147.1]$ & .869 & $2.073\ [1.473,2.722]$ & .902 \\
LLaVA (AV) & $1.23\ [-.25,2.96]$ & .393 & $.001\ [-.005,.007]$ & .459 \\
LLaVA (AA) & $500.2\ [124.9,1{,}017.9]$ & .836 & $1.563\ [1.106,2.072]$ & .869 \\
\bottomrule
\end{tabular*}
\end{minipage}
\end{center}

The independently selected 52-example position grid reproduces the same position-group preference. The mean logit restoration benefit under answer-region corruption exceeds that under same-area-random corruption in all three models, with all-visual positions preferred for Gemma and Qwen and answer-adjacent positions for LLaVA. This replication establishes that the position difference is not specific to the expanded 61-example set.

\subsubsection{Frozen answer-position controls confirm content, direction, and location}

The position-grid result identifies where restoration is clearest, but it does not by itself show that the restored content and direction are specific. We therefore use a separate 60-image development set to choose, for each model, the earliest layer whose mean answer-position restoration reaches at least 90\% of that model's development-set maximum. The selected layers are then frozen and evaluated on the same 240 confirmation images without changing the layer, target answer, four answer-adjacent positions, or control count. Table~\ref{tab:answer-position-confirmation} reports the confirmation results for the restoration itself and the three specificity controls.

\begin{samepage}
\begin{center}
\begin{minipage}{\linewidth}
\captionsetup{type=table,hypcap=false}
\caption{\textbf{Frozen answer-position restoration is specific in all three models.} The Restore column reports the mean increase in complete-answer log probability over the corrupted baseline; the three control columns are paired contrasts between correct restoration and the named intervention, not additional restoration magnitudes. A control contrast can exceed Restore when that control lowers complete-answer log probability below the corrupted baseline. Brackets are 95\% cluster-bootstrap intervals on the same 240 confirmation images. Layers were selected on a separate 60-image development set. All 12 comparisons remain positive after joint correction.}
\label{tab:answer-position-confirmation}
\centering
\small
\setlength{\tabcolsep}{2.4pt}
\begin{tabular*}{\linewidth}{@{\extracolsep{\fill}}lccccc@{}}
\toprule
Model & Layer & Restore & vs. random & vs. reverse & vs. wrong pos. \\
\midrule
Gemma & 28 & \shortstack{$3.956$\\$[2.150,5.781]$} & \shortstack{$3.743$\\$[1.970,5.587]$} & \shortstack{$9.778$\\$[6.960,12.611]$} & \shortstack{$3.135$\\$[1.444,4.924]$} \\
Qwen & 23 & \shortstack{$1.831$\\$[1.086,2.628]$} & \shortstack{$2.108$\\$[1.327,2.949]$} & \shortstack{$5.179$\\$[3.947,6.432]$} & \shortstack{$1.823$\\$[1.082,2.613]$} \\
LLaVA & 20 & \shortstack{$1.060$\\$[0.779,1.359]$} & \shortstack{$1.105$\\$[0.825,1.409]$} & \shortstack{$2.367$\\$[1.833,2.946]$} & \shortstack{$0.987$\\$[0.720,1.274]$} \\
\bottomrule
\end{tabular*}
\end{minipage}
\end{center}
\end{samepage}


Correct restoration is compared with the corrupted baseline, four equal-norm random directions, the reverse of the clean-state difference, and the same difference inserted at the wrong text positions. All four comparisons pass for each model after the 12 tests are corrected together. Randomizing the restored vector, reversing its direction, or inserting it at the wrong positions each reduces the recovery; the consistent reduction ties the restoration gain to the correct internal content, direction, and location. Combined with the regional comparison, the four specificity controls establish the same operational pattern in all three models: answer-region corruption causes greater score loss than matched random corruption, and the aligned clean hidden state selectively recovers that additional loss. The next question is whether a fitted sparse reconstruction preserves the internal change caused by masking the annotated region and produces the corresponding target-score change at its tested feed-forward output.

\FloatBarrier

\subsection{The current sparse tools do not reproduce a positive target-score change}
\label{sec:sparse-effect}

Restoring the complete hidden-state difference increases target-answer scores, but the increase does not show that a fitted sparse reconstruction produces the same score change. We compare complete and sparse interventions at one matched internal object: the feed-forward output at layer 17 for Gemma, layer 14 for Qwen, and layer 16 for LLaVA. For each model, the experiment measures the feed-forward output in the complete-image and answer-region-corrupted runs, reconstructs the output difference with the sparse tool, and writes each candidate difference back at the same position. The combined intervention uses 32 positions sampled uniformly across the visual tokens and the four text positions immediately preceding the answer. Keeping the examples, answer readout, and write locations fixed separates the complete change, sparse reconstruction, and reconstruction residual. We first qualify reconstruction accuracy and then evaluate how inserting or deleting each component changes the complete-answer log probability.

\subsubsection{The sparse tools reconstruct the relevant change poorly}

The fitted sparse tools do not accurately reconstruct the change caused by evidence corruption in the compact-evidence questions studied here. At the primary $K=128$ budget, Equations~\ref{eq:relative-reconstruction-error} and~\ref{eq:reconstruction-direction-similarity} give mean answer-position errors of $1.756$ (Gemma, $n=60$), $1.872$ (Qwen, $n=60$), and $1.617$ (LLaVA, $n=59$), with corresponding direction similarities of $.227$, $.221$, and $.199$. One LLaVA image is excluded from both image-level means because its complete change is zero at all four positions. We therefore interpret the following results as measurements of what each fitted tool reveals, not as direct measurements of a complete sparse mechanism in the base model.

Even under that restricted interpretation, the comparison itself is well controlled. Each model uses the same 240 confirmation images, the same complete-answer log-probability readout, and the same write location for the complete change, sparse reconstruction, and reconstruction residual. At most 128 nonnegative feature activations per token form the main sparse reconstruction; fixed budgets of 16, 32, 64, and 128 are all retained. Original-state writeback, vector decomposition, joint writeback, random-direction norm matching, repeated forward passes, and agreement with the native-channel answer baselines all pass their numerical checks.

\subsubsection{Sparse and remaining components change target-answer scores differently}

Table~\ref{tab:same-object-components} reports the combined visual- and answer-position interventions at the main 128-feature budget. Positive restoration means that inserting a component from the complete-image run increases the correct answer's complete-sequence log probability in the corrupted run. Positive deletion means that removing the corresponding complete--corrupted difference decreases the same sequence score.

\begin{samepage}
\begin{center}
\begin{minipage}{\linewidth}
\captionsetup{type=table,hypcap=false}
\caption{\textbf{The selected sparse reconstruction does not produce a reliable positive target-score change.} Results combine 32 uniformly sampled visual positions with the four text positions immediately preceding the answer and use at most 128 active features per token on the same 240 images per model. Each cell gives the mean change in complete-answer log probability and an uncorrected 95\% cluster-bootstrap interval. Sparse is the selected reconstruction; residual is the reconstruction residual. $\ast$ marks results that remain different from zero after joint correction of 72 tests.}
\label{tab:same-object-components}
\centering
\small
\setlength{\tabcolsep}{2.2pt}
\begin{tabular*}{\linewidth}{@{\extracolsep{\fill}}lcccc@{}}
\toprule
Model & Sparse restore & Residual restore & Sparse delete & Residual delete \\
\midrule
Gemma & \shortstack{$-0.428^{\ast}$\\$[-0.599,-0.260]$} & \shortstack{$0.516^{\ast}$\\$[0.281,0.760]$} & \shortstack{$-0.540^{\ast}$\\$[-0.706,-0.372]$} & \shortstack{$0.314$\\$[0.110,0.515]$} \\
Qwen & \shortstack{$-0.005$\\$[-0.028,0.017]$} & \shortstack{$0.022$\\$[0.002,0.043]$} & \shortstack{$-0.008$\\$[-0.026,0.010]$} & \shortstack{$-0.007$\\$[-0.028,0.012]$} \\
LLaVA & \shortstack{$-0.017$\\$[-0.033,-0.002]$} & \shortstack{$0.000$\\$[-0.017,0.018]$} & \shortstack{$0.003$\\$[-0.008,0.015]$} & \shortstack{$0.010$\\$[-0.006,0.025]$} \\
\bottomrule
\end{tabular*}
\end{minipage}
\end{center}
\end{samepage}


No positive sparse-component score change survives the joint correction in any model. Gemma instead shows negative sparse restoration and deletion changes, while restoring the reconstruction residual increases the sequence score. Qwen and LLaVA show no component-level change after the same correction. The specificity controls reinforce this boundary: no sparse reconstruction outperforms equal-norm random directions or a wrong insertion position, and Gemma's combined-position sparse restoration is significantly weaker than both controls.

The four fixed feature budgets do not reveal a hidden positive pattern. Their mean sparse-restoration changes reverse direction across budgets and remain small for Qwen and LLaVA; Gemma's answer-position change remains negative from 16 through 128 features. These curves are descriptive because the primary joint tests use the frozen 128-feature setting. The conclusion is therefore about the current tools: the tested sparse reconstructions do not produce a reliable positive target-score change under restoration or deletion. This failure does not imply that the base models lack sparse computation, because reconstruction error may prevent the fitted interfaces from reproducing the complete internal change. Before removing the fitted tools from the next causal test, we examine what their sparse descriptions display across models and why those descriptions cannot by themselves establish a path.

\FloatBarrier

\subsection{The sparse tools display different normalized structures across models}
\label{sec:sparse-structure}

The preceding same-object intervention shows that the fitted sparse reconstructions do not reliably recover target-answer scores. Their feature patterns may nevertheless describe how each fitted tool responds to the image corruption. We therefore ask how the clean--corrupted change is distributed within each sparse representation. Because feature identities are not shared across models, we compare normalized summaries rather than raw feature numbers, and we do not use these summaries to infer a faithful path.

\begin{samepage}
The clean--masked support overlap is the fraction of active feature identities shared by the two image conditions. The feature fraction for 80\% of change is the proportion of the tool's full feature set needed to account for 80\% of the squared sparse change; a smaller value means that the change is more concentrated. The positive change-energy fraction measures how much of that squared change comes from features that increase rather than decrease. Finally, the log sparse/full norm ratio compares the size of the sparse reconstruction with the complete internal change within the same model and position group. Table~\ref{tab:sparse-structure-summary} reports these four summaries across models and position groups.
\end{samepage}

\begin{samepage}
\begin{center}
\begin{minipage}{\linewidth}
\captionsetup{type=table,hypcap=false}
\caption{\textbf{The structures seen by the sparse tools differ across models.} Values are model means on the same 240 images. ``Pairs'' counts how many of the three paired model comparisons remain different after the joint correction of 72 tests. These summaries describe each fitted tool and do not align feature identities across models.}
\label{tab:sparse-structure-summary}
\centering
\small
\setlength{\tabcolsep}{2.2pt}
\begin{tabular*}{\linewidth}{@{\extracolsep{\fill}}llrrrr@{}}
\toprule
Summary & Positions & Gemma & Qwen & LLaVA & Pairs \\
\midrule
Clean--masked support overlap & Visual & 0.351 & 0.406 & 0.407 & 2/3 \\
Clean--masked support overlap & Answer & 0.677 & 0.690 & 0.684 & 0/3 \\
Feature fraction for 80\% of change & Visual & 0.0475 & 0.0660 & 0.0345 & 3/3 \\
Feature fraction for 80\% of change & Answer & 0.0142 & 0.0099 & 0.0083 & 3/3 \\
Positive change-energy fraction & Visual & 0.512 & 0.525 & 0.510 & 1/3 \\
Positive change-energy fraction & Answer & 0.491 & 0.495 & 0.519 & 2/3 \\
Log sparse/full norm ratio & Visual & -0.159 & -0.247 & -0.151 & 2/3 \\
Log sparse/full norm ratio & Answer & 0.079 & 0.125 & -0.001 & 3/3 \\
\bottomrule
\end{tabular*}
\end{minipage}
\end{center}
\end{samepage}


Sixteen of the 24 paired model comparisons remain different after all 72 tests are corrected together. The clearest pattern is concentration. At visual positions, the feature fractions needed for 80\% of the change are $.0660$ for Qwen, $.0475$ for Gemma, and $.0345$ for LLaVA; at answer positions, they are $.0142$ for Gemma, $.0099$ for Qwen, and $.0083$ for LLaVA. All six pairwise comparisons for this summary remain different after correction. The other summaries are less uniform: clean--masked support overlap differs for two visual-position pairs but no answer-position pair, while positive change energy differs for one visual-position pair and two answer-position pairs.

These results establish model-conditioned structural differences in what the fitted tools display. They do not establish a shared feature identity, and they do not show that the base models themselves use the same or different complete sparse mechanisms. The earlier intervention found no reliable positive target-score change from these reconstructions despite their different normalized structures. We next remove the fitted tools from the intervention and test whether selected channels belonging to the base models account for the score change.

\FloatBarrier

\subsection{Selected native channels do not form a complete faithful path}
\label{sec:native-path}

The fitted sparse tools show different structures across models, but their small or inconsistent sequence-score changes may reflect the tools rather than the base models. Native channels provide a separate test without fitted reconstruction error. A native channel is one coordinate after the gated feed-forward calculation and before the layer projects that calculation back into the residual stream. It is therefore a pre-projection base-model coordinate, not a coordinate of the residual-space feed-forward output reconstructed in Section~\ref{sec:sparse-effect}. Native channels are not assumed to be individually meaningful features.

The upstream group covers all visual-token positions in the first half of each model's language layers. The downstream group covers the four answer-adjacent text positions in the second half. This split is fixed at half of the language network and is not chosen from answer effects.

On the 60-image development set, each image's known correct answer is used to score and rank its channels. The development results freeze the selection rule and retained fraction, rather than one shared list of channels. One rule retains only channels with the largest positive contributions to the answer. A second rule retains channels with the largest positive or negative contributions, preventing important suppressive contributions from being discarded. The first rule freezes 1\% per layer in all three models; the second freezes 10\% for Gemma, 5\% for Qwen, and 10\% for LLaVA. The same rule and fraction are then applied separately to each confirmation image using that image's known correct answer. The selected confirmation set is therefore answer-conditioned and image-specific. The experiment tests whether such compact per-image support sets satisfy additional path criteria; it does not identify one answer-independent channel set shared across images.

A compact path must satisfy more than a large selected-group sequence-score change. The selected upstream and downstream groups must jointly change the score more than equal-norm random groups; fixing the downstream group in its corrupted state must weaken the upstream restoration; and restoring the downstream group must recover part of an upstream deletion. A faithful path must additionally stay within 80--120\% of the corresponding full-channel score change, preserve the answer after channels outside the selection are removed, and show little recovery from those outside channels alone. Table~\ref{tab:native-channel-path} evaluates both selection rules against these criteria.

\begin{samepage}
\begin{center}
\begin{minipage}{\linewidth}
\captionsetup{type=table,hypcap=false}
\caption{\textbf{Strong native-channel effects do not establish a faithful compact path.} Restore and delete columns report two raw complete-answer log-probability effects separated by a slash: selected channels first, full-channel reference second. The slash does not denote a ratio. Positive-only selection retains 1\% of channels per layer; the second selection retains the largest positive and negative contributions. ``Local link'' requires the upstream blocking and downstream bypass tests to beat matched random groups after the 66-comparison correction. A faithful path additionally requires effect matching and the tests of channels outside the selection.}
\label{tab:native-channel-path}
\centering
\small
\setlength{\tabcolsep}{1.8pt}
\begin{tabular*}{\linewidth}{@{\extracolsep{\fill}}llccccc@{}}
\toprule
Model & Selection & Kept & Restore selected/full & Delete selected/full & Local link & Faithful path \\
\midrule
Gemma & Positive only & 1\% & 16.472/3.791 & 50.084/3.294 & Yes & No \\
Gemma & Positive + negative & 10\% & 3.940/3.791 & 3.237/3.294 & No & No \\
Qwen & Positive only & 1\% & 4.506/1.859 & 12.709/1.706 & Yes & No \\
Qwen & Positive + negative & 5\% & 1.869/1.859 & 1.414/1.706 & Yes & No \\
LLaVA & Positive only & 1\% & 2.972/-0.288 & 7.082/1.935 & Yes & No \\
LLaVA & Positive + negative & 10\% & 0.474/-0.288 & 0.815/1.935 & No & No \\
\bottomrule
\end{tabular*}
\end{minipage}
\end{center}
\end{samepage}


Positive-only selection produces a local upstream--downstream link in all three models, but its effects are not faithful to the full-channel reference. Gemma restoration is $16.472$ compared with a full-channel effect of $3.791$, and its deletion effect is $50.084$ compared with $3.294$. Qwen and LLaVA show the same over-amplification pattern. Selecting both positive and negative contributions brings the Gemma and Qwen restoration means close to their full-channel references, but the full set of criteria still fails.

The positive-only local links survive the 66-comparison correction in all three models, but their strong over-amplification fails full-channel effect matching. Under the positive-and-negative rule, only Qwen shows an answer-conditioned local upstream--downstream link: its blocking and bypass advantages over matched random groups have corrected $p$-values of $.0138$ and $.0438$. Qwen still does not pass the complete effect-matching and outside-channel tests. Gemma does not show a specific downstream handoff under this rule, and LLaVA's full-channel restoration reference is itself negative and uncertain. No model therefore provides a complete, faithful native-channel path under either frozen rule.

The native-channel test removes the fitted sparse tool from the intervention, yet no selected channel group forms a complete faithful path. This rules out the tested axis-aligned groups, not a distributed combination of channels. The remaining compact alternative is therefore a non-axis-aligned subspace shared across examples. The next section tests that alternative before returning to the complete hidden-state change.

\FloatBarrier

\subsection{Shared low-rank candidates fail; one LLaVA full-hidden interface passes}
\label{sec:subspace-bottleneck}

The sparse and native-channel tests leave a specific alternative unresolved. The clean-minus-masked hidden-state change may be distributed across individual coordinates while remaining confined to a low-dimensional residual-state subspace shared across examples within each model. We test this possibility before interpreting the complete hidden-state intervention. The experiment fits a separate basis for each model at frozen interfaces: layer 1 at visual positions for Gemma, layer 14 at visual positions for Qwen, and layer 30 at the four answer-adjacent positions for LLaVA.

\subsubsection{No example-shared rank from 1 to 128 passes the bidirectional gate}

For each model, 40 images fit the uncentered basis in Equation~\ref{eq:shared-subspace-projection}; 20 different images select among ranks $1,2,4,8,16,32,64,$ and $128$; and a separate 40-image set is reserved for replication. A candidate must restore and remove between 80\% and 120\% of the mean clean-minus-masked sequence-score change on the selection images. No tested rank passes both mean gates in any model, so no rank advances to strict confirmation. Table~\ref{tab:shared-subspace-summary} reports the largest tested rank on the replication images together with the full-hidden reference at the same interface.

\begin{samepage}
\begin{center}
\begin{minipage}{\linewidth}
\captionsetup{type=table,hypcap=false}
\caption{\textbf{An example-shared 128-dimensional subspace fails bidirectional effect matching at the initial interfaces.} Gemma/Qwen use visual positions and LLaVA uses answer-adjacent positions, each at its initially selected layer. Values are percentages of the mean clean-minus-masked sequence-score change on the separate 40-example replication set. Energy is the squared hidden-difference norm retained by the projection. Values above 100 indicate effect overshoot, and negative values indicate effect reversal. Rank 128 is the largest tested subspace; no rank from 1 to 128 passes both restoration and corrupted-state-replacement selection gates in any model. Full columns report the unprojected hidden-state reference at the same layer and positions.}
\label{tab:shared-subspace-summary}
\centering
\normalsize
\setlength{\tabcolsep}{2.8pt}
\begin{tabular*}{\linewidth}{@{\extracolsep{\fill}}lrrrrr@{}}
\toprule
Model & Rank-128 energy & Rank-128 restore & Rank-128 delete & Full restore & Full delete \\
\midrule
Gemma & 88.9 & 425.0 & -357.7 & 192.3 & 83.2 \\
Qwen & 53.1 & 53.0 & 51.2 & 94.6 & 76.5 \\
LLaVA & 15.9 & 10.5 & 0.1 & 94.9 & 92.4 \\
\bottomrule
\end{tabular*}
\end{minipage}
\end{center}
\end{samepage}


Retaining substantial hidden-difference energy is not sufficient for causal fidelity. Gemma's rank-128 projection retains 88.9\% of the squared difference but over-restores the answer effect and reverses the deletion effect. Qwen retains 53.1\% of the energy and approximately half of both causal effects. LLaVA retains 15.9\% of the energy, restores 10.5\% of the answer effect, and removes 0.1\%. Failure at every selection rank therefore rejects an example-shared 1--128-dimensional complete subspace at these frozen interfaces. The tested rank range does not rule out higher-rank, nonlinear, or sample-specific representations.

\subsubsection{A result-blind evaluation confirms the LLaVA full-hidden interface}

At the initial frozen interfaces, the unprojected reference gives different results. Gemma's layer-1 visual-position restoration is unstable across the selection and replication sets, while Qwen's layer-14 visual-position replacement effect remains below the 80\% gate. The evaluated LLaVA layer-30 answer-position interface keeps both mean effects within 80--120\% on both sets: restoration and replacement are 93.0\% and 91.0\% on selection and 94.9\% and 92.4\% on replication. We therefore freeze only this LLaVA interface for a separate result-blind test; no layer, position, rank, or threshold is rescanned. Because the three initial interfaces differ in both layer and position type, their cross-model contrast alone cannot assign the difference to model architecture.

The confirmation pool contains 61 masks reviewed without model predictions or intervention results. Human review accepts 60 masks and rejects one before model evaluation. On the 60 accepted examples, the mean clean-minus-masked sequence-score change is $.9378$. Restoring the complete hidden-state difference yields $.8981$, or 95.8\% of that mean change; replacing the clean hidden states with their masked values removes $.8907$, or 95.0\% of the same reference. Percentages use the unrounded means. Table~\ref{tab:llava-hidden-confirmation} shows that all four predeclared paired-bootstrap gates have positive lower bounds.

\begin{samepage}
\begin{center}
\begin{minipage}{\linewidth}
\captionsetup{type=table,hypcap=false}
\caption{\textbf{The frozen LLaVA full-hidden interface passes separate result-blind confirmation.} The interface is layer 30 at the four text positions immediately preceding the answer. Restore and Delete are percentages of the mean clean-minus-masked sequence-score change; Delete writes masked hidden states into the clean run rather than zeroing them. R-L and R-U are raw sequence-log-probability lower bounds for restoration minus 80\% of that mean change and 120\% of that mean change minus restoration; D-L and D-U are the corresponding replacement bounds. All four bounds must be nonnegative, and Gate 1 denotes a pass.}
\label{tab:llava-hidden-confirmation}
\centering
\normalsize
\setlength{\tabcolsep}{2.2pt}
\begin{tabular*}{\linewidth}{@{\extracolsep{\fill}}lrrrrrrrr@{}}
\toprule
Interface & $n$ & Restore & Delete & R-L & R-U & D-L & D-U & Gate \\
\midrule
LLaVA L30, AA-4 & 60 & 95.8 & 95.0 & .050 & .097 & .047 & .100 & 1 \\
\bottomrule
\end{tabular*}
\end{minipage}
\end{center}
\end{samepage}


Restoring the evaluated LLaVA full-hidden interface recovers most of the clean-minus-masked sequence-score difference; replacing its clean value with the masked value removes most of that difference. The 60-example result therefore confirms this interface under both specified interventions. It neither establishes a minimum dimension nor identifies an analogous interface in Gemma or Qwen. The initial comparison's layer-and-position asymmetry motivates one further test before we interpret the cross-model pattern.

\subsubsection{A same-sample late-answer comparison narrows the interface asymmetry}

For this follow-up, each model is tested at its penultimate language layer and the four text positions before the answer, using the same 60 reviewed examples and the same complete-answer score. No result is used to rescan the rule within a model. Table~\ref{tab:matched-late-hidden} reports restoration and corrupted-state replacement on the same image--question IDs. The relative-layer rule was motivated by the already observed LLaVA result, so this comparison is diagnostic; the original LLaVA result above remains its independent confirmation.

\begin{samepage}
\begin{center}
\begin{minipage}{\linewidth}
\captionsetup{type=table,hypcap=false}
\caption{\textbf{Only the evaluated LLaVA late-answer interface passes the matched bidirectional test.} All three models use the same 60 reviewed image--question examples, the zero-based penultimate language layer, four answer-adjacent positions, the same complete-answer score, and the same restored/replaced-state protocol in one runtime. $n_+$ counts examples with a positive clean-minus-masked sequence-score change. $V$ is the mean of that change; $R/V$ and $C/V$ are ratios of within-model means in percent for restoration and corrupted-state replacement. R-L and C-L are raw sequence-score 95\% bootstrap lower bounds for the respective effects minus $0.8V$; the upper bounds must also be nonnegative for Gate=1. Negative ratios denote reversed mean effects. Different models' raw score magnitudes are not compared.}
\label{tab:matched-late-hidden}
\centering
\small
\setlength{\tabcolsep}{2.2pt}
\begin{tabular*}{\linewidth}{@{\extracolsep{\fill}}lrrrrrrrr@{}}
\toprule
Model & Layer & $n_+$ & $V$ & $R/V$ & $C/V$ & R-L & C-L & Gate \\
\midrule
Gemma & 32 & 29/60 & 0.407 & -85.0 & -152.8 & -1.837 & -2.196 & 0 \\
Qwen & 26 & 36/60 & 1.463 & 74.3 & 45.3 & -0.483 & -1.017 & 0 \\
LLaVA & 30 & 41/60 & 0.930 & 95.2 & 95.7 & 0.046 & 0.050 & 1 \\
\bottomrule
\end{tabular*}
\end{minipage}
\end{center}
\end{samepage}


Only the evaluated LLaVA interface passes all four paired bounds in this matched run. Gemma's two mean changes reverse direction; Qwen restores 74.3\% but removes only 45.3\% of its mean clean-minus-masked sequence-score change under corrupted-state replacement. The corresponding LLaVA ratios are 95.2\% and 95.7\%. All 60 examples are scored in every model, but the clean-minus-masked sequence-score change is positive in 29, 36, and 41 examples for Gemma, Qwen, and LLaVA, respectively. Thus, the matched result concerns these frozen interfaces on this LLaVA-origin sample set, not the absence of another successful interface in either Gemma or Qwen. The original LLaVA confirmation and this unified-runtime rerun give slightly different point estimates but the same four-gate decision.

The same late-answer rule is also applied to the 40/20/40 shared-subspace splits. At the matched interfaces, none of the eight ranks from 1 to 128 passes both selection means in any model (Appendix~\ref{app:shared-subspace-details}). This removes the initial visual-versus-answer position difference as an explanation for the tested low-rank failure, while leaving higher-rank and sample-specific objects open. Across the evaluated interfaces, all three models show additional damage from masking evidence rather than matched control content; the tested compact candidates fail their own complete-effect criteria, while one full LLaVA hidden-state interface passes independent bidirectional confirmation against the clean-minus-masked sequence-score change. The Discussion considers what can and cannot be inferred from that separation.

\FloatBarrier


\subsection{Sensitivity to the effect-matching tolerance}
\label{app:gate-tolerance}

The primary complete-effect rule uses the frozen 80--120\% mean interval. As a post-hoc robustness check, we recompute the same four paired sample-bootstrap margins on the same 60 independently reviewed LLaVA examples at symmetric tolerances of 10, 15, 20, 25, and 30\%. The 20\% bootstrap lower bounds reproduce the original confirmation exactly. The low-rank column below rechecks only the two selection-set means for the eight ranks in each model; the Qwen column rechecks only the two means of the later matched-interface diagnostic. These auxiliary columns do not constitute independent confirmations.

\begin{center}
\begin{minipage}{\linewidth}
\captionsetup{type=table,hypcap=false}
\caption{\textbf{Bidirectional decisions across tolerance widths.} The LLaVA minimum is the smallest of four 95\% paired-bootstrap lower bounds on the independent 60-example set, in sequence-log-probability units. The low-rank count is the number of 24 model--rank candidates passing both selection means; Qwen is the later same-sample, same-relative-layer mean decision. The latter two columns do not use the independent LLaVA bootstrap gate.}
\label{tab:gate-tolerance}
\centering
\begin{tabular*}{\linewidth}{@{\extracolsep{\fill}}ccccc@{}}
\toprule
Tol. & LLaVA min. bound & LLaVA 4/4 & Low-rank / 24 & Qwen both means \\
\midrule
10\% & $-.023$ & 0 & 0 & 0 \\
15\% & $+.015$ & 1 & 0 & 0 \\
20\% & $+.047$ & 1 & 0 & 0 \\
25\% & $+.076$ & 1 & 0 & 0 \\
30\% & $+.102$ & 1 & 0 & 0 \\
\bottomrule
\end{tabular*}
\end{minipage}
\end{center}

The independent LLaVA point estimates satisfy the 10\% mean gate, but the tighter paired-bootstrap bounds cross zero. Its four-bound decision is positive from 15\% through 30\%. None of the 24 low-rank selection candidates passes both mean gates in this range, and the matched Qwen interface remains outside the mean gate because masked-state writeback retains only 45.3\% of its natural score change. This analysis checks the sensitivity of the original decision; it does not replace the predeclared 20\% criterion.
